\subsubsection{例子}\label{subsec:distribution-derivative-examples}

\begin{Example}
\label{ex:heaviside-derivative}
考虑函数
\begin{equation}\label{eq:heaviside-function-definition}
H(x)=\begin{cases}0,&x\leq0\\1,&x>0.\end{cases}
\end{equation}
计算得到
\begin{equation}\label{eq:heaviside-derivative-delta-pairing}
(H',\phi)=-(H,\phi')=-\int_0^\infty\phi'(x)\,\mathrm{d}x=\phi(0)=(\delta,\phi),
\end{equation}
即 $H$ 的导数为 delta 函数. 函数 $H$ 称为\textbf{Heaviside 函数}.
\end{Example}

\begin{Example}
\label{ex:higher-delta-derivative}
delta 函数的 $k$ 阶导数为泛函 $\phi\mapsto(-1)^k\phi^{(k)}(0)$.
\end{Example}

\begin{Example}
\label{ex:singular-power-derivative}
令
\begin{equation}\label{eq:positive-part-singular-power}
x_+^\lambda=\begin{cases}0,&x\leq0\\x^\lambda,&x>0\end{cases}
\end{equation}
且 $-1<\lambda<0$. 直观上, 可能预期导数为 $\lambda x_+^{\lambda-1}$, 但这一函数具有不可积奇性, 因而不是分布. 正确的答案是作适当的正则化. 计算得到\[
\begin{aligned}
((x_+^\lambda)',\phi)&=-(x_+^\lambda,\phi')\\
&=-\int_0^\infty x^\lambda\phi'(x)\,\mathrm{d}x\\
&=-\lim_{\epsilon\to0}\int_\epsilon^\infty x^\lambda\phi'(x)\,\mathrm{d}x\\
&=\lim_{\epsilon\to0}\int_\epsilon^\infty\lambda x^{\lambda-1}(\phi(x)-\phi(\epsilon))\,\mathrm{d}x\\
&=\int_0^\infty\lambda x^{\lambda-1}(\phi(x)-\phi(0))\,\mathrm{d}x.
\end{aligned}
\]
\end{Example}

\begin{Example}
\label{ex:boundary-jump-derivative}
设 $\Omega$ 为具有光滑边界 $\Gamma$ 的区域. 设 $f$ 属于 $C^1(\overline\Omega)$, 并令在 $\Omega$ 外部 $f=0$. 将 $f$ 视为 $\mathbb{R}^m$ 上的分布. 得到\[
\begin{aligned}
\left(\dfrac{\partial f}{\partial x_j},\phi\right)&=-\left(f,\dfrac{\partial\phi}{\partial x_j}\right)\\
&=-\int_\Omega f(\bm{x})\dfrac{\partial\phi}{\partial x_j}(\bm{x})\,\mathrm{d}\bm{x}\\
&=\int_\Omega\dfrac{\partial f}{\partial x_j}(\bm{x})\phi(\bm{x})\,\mathrm{d}\bm{x}
-\int_\Gamma f(\bm{x})\phi(\bm{x})n_j\,\mathrm{d}S.
\end{aligned}
\]这里 $n_j$ 为 $\Gamma$ 的单位外法向的第 $j$ 个分量, $\mathrm{d}S$ 为 $m-1$ 维曲面面积微元. 因此, $f$ 的分布导数包含两项, 一项对应于 $\Omega$ 内的普通导数, 另一项涉及支集在 $\Gamma$ 上的分布. 后一项来自 $f$ 跨越 $\Gamma$ 时的跳跃.
\end{Example}

\begin{Example}
\label{ex:polarization-charge}
前一个例子在电磁学中有一些重要应用. 设 $\Omega\subset\mathbb{R}^3$ 为具有光滑边界 $\Gamma$ 的区域. 假设有一个属于 $C^1(\overline\Omega)$ 的\textbf{极化}向量场 $\bm{p}:\overline\Omega\to\mathbb{R}^3$. 令在 $\Omega$ 外部 $\bm{p}=0$, 就可以将 $\bm{p}$ 视为 $\mathbb{R}^3$ 上的分布. 于是将\textbf{极化电荷}定义为 $\bm{p}$ 在\emph{分布意义下}的散度. 计算如下.\[
\begin{aligned}
(\nabla\cdot\bm{p},\phi)&=-\sum_{i=1}^3\left(p_i,\dfrac{\partial\phi}{\partial x_i}\right)\\
&=-\int_\Omega\bm{p}(\bm{x})\cdot\nabla\phi(\bm{x})\,\mathrm{d}\bm{x}\\
&=\int_\Omega\nabla\cdot\bm{p}(\bm{x})\phi(\bm{x})\,\mathrm{d}\bm{x}
-\int_\Gamma\bm{p}(\bm{x})\cdot\bm{n}(\bm{x})\phi(\bm{x})\,\mathrm{d}A.
\end{aligned}
\]这里 $\bm{n}$ 为 $\Gamma$ 的单位外法向, $\mathrm{d}A$ 为 $\Gamma$ 上的曲面面积微元. 因此, 极化电荷包含一项对应于 $\Omega$ 内 $\bm{p}$ 的普通散度, 以及由 $\Gamma$ 上 $\bm{p}$ 的法向分量给出的面电荷. 后一项来自 $\bm{p}$ 跨越 $\Gamma$ 时的跳跃. 若 $\bm{p}$ 分片光滑, 并在 $\Omega$ 内部具有跳跃间断曲面, 则沿这些曲面的跳跃的法向分量也会贡献极化电荷.
\end{Example}

\begin{Example}
\label{ex:inverse-radius-laplacian}
在 $\mathbb{R}^3$ 上令 $f(\bm{x})=1/|\bm{x}|=1/r$. 容易验证, 当 $r\ne0$ 时 $\Delta(1/r)=0$. 将在分布意义下计算 $1/r$ 的 Laplacian. 计算得到\[
\begin{aligned}
\left(\Delta\dfrac1r,\phi\right)&=\left(\dfrac1r,\Delta\phi\right)\\
&=\int_{\mathbb{R}^3}\dfrac{\Delta\phi(\bm{x})}r\,\mathrm{d}\bm{x}\\
&=\lim_{\epsilon\to0}\int_{r\geq\epsilon}\dfrac{\Delta\phi}r\,\mathrm{d}\bm{x}.
\end{aligned}
\]分部积分给出
\begin{equation}\label{eq:inverse-radius-laplacian-integration-by-parts}
\begin{aligned}
\int_{r\geq\epsilon}\dfrac{\Delta\phi}r\,\mathrm{d}\bm{x}
&=\int_{r\geq\epsilon}\Delta\left(\dfrac1r\right)\phi\,\mathrm{d}\bm{x}
-\int_{r=\epsilon}\dfrac{\partial\phi}{\partial r}\dfrac1r\,\mathrm{d}S\\
&\quad+\int_{r=\epsilon}\phi\dfrac{\partial}{\partial r}\dfrac1r\,\mathrm{d}S.
\end{aligned}
\end{equation}
在 \eqref{eq:inverse-radius-laplacian-integration-by-parts} 的右端, 第一项为零, 当 $\epsilon\to0$ 时第二项为 $\epsilon$ 阶, 最后一项等于 $-\epsilon^{-2}\int_{r=\epsilon}\phi\,\mathrm{d}S$, 即半径为 $\epsilon$ 的球面上 $\phi$ 的平均值的 $-4\pi$ 倍. 因而令 $\epsilon\to0$, 得到
\begin{equation}\label{eq:inverse-radius-laplacian-delta-pairing}
\left(\Delta\dfrac1r,\phi\right)=-4\pi\phi(0),
\end{equation}
即
\begin{equation}\label{eq:inverse-radius-laplacian-delta}
\Delta\dfrac1r=-4\pi\delta.
\end{equation}
方程 $Lu=\delta$ 的解相当重要, 其中 $L$ 为常系数偏微分算子; 下两节将进一步研究这类解.
\end{Example}
